\begin{frame}{What this design does}

\begin{exampleblock}{\faCompassDrafting~Transferable principles}
For CDIO-aligned, industry-engaged courses
\end{exampleblock}

\begin{itemize}
    \item sustained industry context (no repeated onboarding)
    \item assessment drives iteration (not just evaluation)
    \item contributions are visible and reviewable
    \item feedback accumulates across the course
    \item students adapt to different audiences
    \item tooling supports -- but does not define -- the design
\end{itemize}

\end{frame}
\note{
This can be summarized in a few key design choices.  
\begin{enumerate}
    \item First, a sustained industry context --  
so students build depth instead of repeatedly onboarding.  
    \item Second, assessment is used to drive iteration,  
not just to evaluate final outputs.  
    \item Third, contributions are made visible and reviewable,  
which supports accountability without fragmenting teamwork.  
    \item Fourth, feedback accumulates across the course,  
so students revisit and improve earlier work.  
    \item And finally, students learn to adapt their work  
to different audiences -- academic and industry.  

\end{enumerate}

The tooling supports this,  
but the key contribution is the alignment  
between learning and professional practice.
}